\documentclass[11pt,a4paper]{article}
\usepackage[margin=25mm]{geometry}
\usepackage{amsmath,amssymb,bm,mathtools}
\usepackage[T1]{fontenc}
\usepackage{lmodern,microtype}
\usepackage[colorlinks=true,linkcolor=blue,urlcolor=blue]{hyperref}
\newcommand{\R}{\mathbb R}
\newcommand{\E}{\mathbb E}
\title{The missing $1/N$ in the hidden-layer kernel $C$\\(Section 2, linear one-hidden-layer case)}
\author{}
\date{}
\begin{document}
\maketitle

\section{Problem}
The uncorrected text of Section~2 writes the hidden-layer kernel of the one-hidden-layer linear network as
\begin{equation}
  \Phi^{\mu,\nu}_{\alpha,\beta} = k_1 \delta_{\alpha,\beta}\, C^{\mu,\nu},
  \qquad
  C^{\mu,\nu} = \big\langle \phi^\mu_i K_{i,j} \phi^\nu_j \big\rangle
  \quad\text{(sum over } i,j = 1,\dots,N\text{)},
  \label{eq:C-as-written}
\end{equation}
and consequently writes $C = \sum_q \lambda_q \Sigma_q$ and $D_C = \sum_q \lambda_q D_q + \delta I_P$.
The definition of the kernels in Section~1 (equation \texttt{eq:kernel-def}) is
\begin{equation}
  \Phi^{(l)\mu,\nu}_{i,j} = \frac1N \sum_{a,b=1}^N \phi^{(l)\mu}_a K^{(l)}_{i,a,j,b}\, \phi^{(l)\nu}_b ,
\end{equation}
so the correct hidden-layer kernel carries a factor $1/N$:
\begin{equation}
  C^{\mu,\nu} = \frac1N \sum_{i,j} K_{i,j} \big\langle \phi^\mu_i \phi^\nu_j \big\rangle
  = \frac1N \sum_q \lambda_q \Sigma^{\mu,\nu}_q ,
  \qquad
  D_C = \frac1N \sum_q \lambda_q D_q + \delta I_P ,
  \qquad \delta = \frac{1}{k_1\kappa}.
  \label{eq:C-correct}
\end{equation}

\section{Derivation from first principles}
The $1/N$ is not a convention that can be chosen freely; it follows from the network.
The readout is $f^\mu_\alpha = N^{-1}\sum_j W^{(1)}_{\alpha,j}\phi^\mu_j$ (mean-field scaling, $\sigma^{(L)} = N$), with Gaussian prior
$\E[W_{\alpha,i} W_{\beta,j}] = K^{(1)}_{\alpha,i,\beta,j} = k_1 \delta_{\alpha,\beta} K_{i,j}$.
Integrating out $W^{(1)}$ against the source $\frac{i}{N}\hat f^\mu_\alpha \phi^\mu_i$ gives
\begin{equation}
  \exp\Big(-\frac{1}{2N^2}\,\hat f^\mu_\alpha \phi^\mu_i K^{(1)}_{\alpha,i,\beta,j}\phi^\nu_j \hat f^\nu_\beta\Big)
  = \exp\Big(-\frac{1}{2N}\,\hat f^\mu_\alpha \Phi^{\mu,\nu}_{\alpha,\beta}\hat f^\nu_\beta\Big),
  \qquad
  \Phi = \frac1N\, \phi K^{(1)} \phi ,
\end{equation}
and it is precisely this $\Phi$ (with its $1/N$) that appears in the output term
$-\frac N2\, y\,(\Phi + \kappa^{-1}I)^{-1} y$ of the action after integrating out $\hat f$ and $f$.
The constraint is imposed by
$\int \mathcal D\hat\Phi\, \exp\{-\tfrac{iN}{2}\hat\Phi\,[\Phi - \tfrac1N \phi K^{(1)}\phi]\}$, so the stationarity condition in $\hat\Phi$ is
\begin{equation}
  \Phi^{\mu,\nu}_{\alpha,\beta} = \frac1N \big\langle \phi^\mu_i K^{(1)}_{\alpha,i,\beta,j}\phi^\nu_j\big\rangle
  = k_1\delta_{\alpha,\beta}\,\frac1N \sum_{i,j}K_{i,j}\langle\phi^\mu_i\phi^\nu_j\rangle .
\end{equation}
The $N$ in the constraint and the $1/N$ in the kernel cancel in the \emph{single-site} tilt
($\frac{i}{2}\phi\,\hat\Phi K^{(1)}\phi$), which is why the tilt $B K$ in Section~2 is correctly written without an $N$;
but they do \emph{not} cancel in the self-consistency for $C$.
In the Fourier basis, with $\langle \phi^\mu_i\phi^\nu_j\rangle = \sum_q F_{i,q}\Sigma_q^{\mu,\nu}F^*_{j,q}$ and $F$ unitary,
$\sum_{i,j}K_{i,j}\langle\phi_i\phi_j\rangle = \operatorname{tr}(K F\Sigma F^\dagger) = \sum_q\lambda_q\Sigma_q$, which gives \eqref{eq:C-correct}.

The factor cannot be absorbed into $k_1$ alone. Write $C' = \sum_q\lambda_q\Sigma_q = NC$ for the uncorrected kernel. Then
$\tilde C' = C' + \frac{1}{k_1\kappa} = N\big(C + \frac{1}{Nk_1\kappa}\big)$ and
$B' = \frac{1}{k_1}\tilde C'^{-1}YY^\top\tilde C'^{-1} = \frac{1}{N^2k_1}\big(C+\frac{1}{Nk_1\kappa}\big)^{-1}YY^\top\big(C+\frac{1}{Nk_1\kappa}\big)^{-1}$.
This is the correct system only after the two substitutions $k_1 \to Nk_1$ \emph{and} $YY^\top \to YY^\top/N$.
So the uncorrected equations describe a network whose readout prior is $N$ times larger and whose targets are $\sqrt N$ times smaller than the ones actually specified.
Nor can it be absorbed into the spatial kernel. Even if $K^{(1),\text{pre}}$ were independent of $K^{(0),\text{post}}$, rescaling $K^{(1),\text{pre}} \to K^{(1),\text{pre}}/N$ to turn $\langle\phi K^{(1),\text{pre}}\phi\rangle$ into the correct $\frac1N\langle\phi K^{(1),\text{pre}}\phi\rangle$ would also rescale the single-site tilt $\frac12\phi\,(B K^{(1),\text{pre}})\phi$ by $1/N$, because the tilt is built from the same matrix; the factor $N$ reappears there.

\section{Consequences}
\paragraph{Large-$N$ scaling.}
The eigenvalues $\lambda_q$ of a translation-invariant kernel of fixed range are $O(1)$, and there are $N$ of them.
\begin{itemize}
  \item With the $1/N$, $C$ is the spectral \emph{average} $\frac1N\sum_q \lambda_q \Sigma_q$, which is $O(1)$ and has a continuum limit
  $\frac1N\sum_q \to \int_{\mathrm{BZ}} \frac{\mathrm d^d q}{(2\pi)^d}$ (lattice units) as $N\to\infty$ at fixed lattice spacing.
  The diagonal element $c$ of $D_C$ is $O(1)$ and the target tilt $b_q \propto \lambda_q \hat s_Y^2 / c^2$ stays finite: the hidden layer learns.
  \item Without it, $c = \sum_q \lambda_q d_q + \delta = O(N)$. Then the target part of the tilt, $\lambda_q\hat s_Y^2/c^2$, vanishes as $N\to\infty$ and only the rate penalty survives, $b_q \to -r$. Hence $d_q \to (a_q^{-1}+r)^{-1} = h_q/\lambda_q$ (the prior covariance shrunk by the rate penalty), and the self-consistency degenerates to $c \simeq \sum_q h_q + \delta$, the $\hat s_Y = 0$ solution of the companion note: the representation reverts to the prior and all feature learning disappears in the thermodynamic limit, unless $\hat s_Y^2$ is made to grow with $N$.
\end{itemize}
So the missing factor changes the qualitative large-$N$ behaviour of \texttt{eq:linear-soln}, not only a constant.

\paragraph{The implicit solution.}
Combining \eqref{eq:C-correct} with the corrected diagonal element of $D_{B,q}$ (it is $\lambda_q\hat s_Y^2/c^2$, not $\lambda_q\hat s_Y^2/c$; see the companion note \texttt{rg\_diagonal\_reduction.tex}), one has per data mode
\begin{equation}
  d_q = \frac{c^2 h_q/\lambda_q}{c^2 - \hat s_Y^2 h_q},
  \qquad
  c = \frac1N\sum_q \lambda_q d_q + \delta ,
\end{equation}
and hence the corrected form of \texttt{eq:linear-soln},
\begin{equation}
  1 - \frac{\delta}{c} = \frac{c}{N}\sum_q \frac{h_q}{c^2 - \hat s_Y^2 h_q},
  \qquad
  h_q = \frac{\lambda_q(\lambda_q\hat s_X^2 + \sigma^2)}{1 + r(\lambda_q \hat s_X^2 + \sigma^2)},
  \qquad
  \hat s_X^2 = \frac{k_0 s_X^2}{m},\ \hat s_Y^2 = \frac{s_Y^2}{k_1}.
  \label{eq:linear-soln-corrected}
\end{equation}
In the continuum limit the right-hand side becomes $c\int_{\mathrm{BZ}}\frac{\mathrm d^dq}{(2\pi)^d}\,\frac{h(q)}{c^2-\hat s_Y^2 h(q)}$.

\paragraph{Check.}
Every parameter is set in \path{latex/notes/scripts/check_rg_linear_soln.py} (run from the repository root with \path{uv run python latex/notes/scripts/check_rg_linear_soln.py}; a few seconds):
a ring of $N=12$ sites with $K_{i,j} = \exp(-r_{i,j}^2/4)$, $r_{i,j}$ the ring distance (a Gaussian of width $\sqrt2$);
$P=4$ patterns, $m=6$ inputs, $n=5$ outputs and $3$ latents;
$X = ZW_X^\top$, $Y = ZW_Y^\top$ with $Z_{\mu,k}\sim\mathcal U[0,1]$ i.i.d., $W_X$ the $Q$ factor of a $6\times3$ standard normal matrix and $W_Y$ $0.6$ times the $Q$ factor of a $5\times3$ one (drawn in the order $Z, W_X, W_Y$ from \path{numpy.random.default_rng(0)}), so that $XX^\top$ and $YY^\top$ commute;
$k_0 = 0.7$, $k_1 = 1.3$, $\sigma^2 = 0.2$, $r = 0.3$, $\delta = 0.5$;
damped iteration (damping $\tfrac12$, $4000$ steps) of the full $NP\times NP$ map $C \mapsto \frac1N\sum_{i,j}K_{i,j}\langle\phi_i\phi_j\rangle$, started from $C = 5I$, with no ansatz and no scalar reduction.
The eigenvalues of $C+\delta I$ at the fixed point are $\{0.68867925, 0.70201000, 1.10586544, 2.46022252\}$, and the roots of \eqref{eq:linear-soln-corrected} agree to all ten printed digits; the off-diagonal part of $U^\top\tilde C U$ is $10^{-15}$.
The equation as written in the uncorrected text gives $\{2.7642, 2.8840, 5.1682, 9.6999\}$, larger by factors $3.9$ to $4.7$.

\section{Recommended fix}
Insert the $1/N$ in the definitions of $\Phi$ and $C$, in $C = \frac1N\sum_q\lambda_q\Sigma_q$, in $\tilde C$, in $D_C$ and in the self-consistency for $c$ (all marked in red in Section~2), and use \eqref{eq:linear-soln-corrected}.
Any analysis built on \texttt{eq:linear-soln}, in particular the limiting-case analysis and the continuum-limit translation, should start from \eqref{eq:linear-soln-corrected}, where $\frac1N\sum_q$ is the natural spectral average.
\end{document}
